% ============================================================
% sec/02metodologia.tex
%
% C´omo buscaron los cinco papers: bases de datos, palabras clave, criterios de inclusi´on/exclusi´on.
% ============================================================
\section{Metodología de búsqueda}
\label{sec:metodologia}

La búsqueda se realizó en Consensus, complementada con
Connected Papers para explorar el grafo de citación y
localizar trabajos adicionales relevantes, y verificada contra
IEEE Xplore para confirmar autoría, medio de publicación y año de
publicación.

\textbf{Palabras clave:} facial recognition, demographic bias,
demographic differentials, fairness, false match rate (FMR), false
non-match rate (FNMR), presentation attack, face morphing, spoofing,
live facial recognition, surveillance, policing, privacy,
accountability, adversarial attacks.

\textbf{Criterios de inclusión:} Medios de publicación con revisión por pares (IEEE,
ACM, OUP, Springer) o informes técnicos con autoría y fecha
identificables (NIST); enfoque en sesgo algorítmico, uso policial o
robustez del reconocimiento facial; publicados entre 2019 y 2026;
con al menos un hallazgo o métrica reportada.

\textbf{Criterios de exclusión:} Blogs sin autoría identificable,
foros, Wikipedia (solo orientación inicial) y notas periodísticas sin
fuente primaria verificable.

De este proceso resultaron las cinco fuentes principales del informe:
cuatro localizadas mediante Consensus y verificadas en IEEE Xplore, y
una obtenida por trazado de referencias en Connected Papers.